% §2 Background and Related Work
\emph{References note:} all related-work references in this paper were verified against the
arXiv API and Crossref on 2026-09-01, including a 2025--26 sweep that surfaced the
full-duplex and interruption benchmark wave discussed in Section~\ref{sec:background} (\S2.7).

We position this paper against five bodies of work.
For each we state what it contributes and, explicitly, where it stops short of a formal
cessation floor.

\subsection{Incremental processing and revocable output}
The closest formal ancestor of a cessation floor is the incremental dialogue-processing
literature.
Schlangen and Skantze \citep{schlangen2011general} model dialogue systems as incrementally
communicating modules whose output units are produced left to right and can be revised or
cancelled.
Toolkits such as InproTK instantiate precisely this revocability for incremental speech
synthesis \citep{baumann2012prosodic}.
That line of work established that output increments are \emph{cancellable} behaviours, but it
was developed for pre-LLM statistical components, carries no agent-orchestration or
tool-execution semantics, and offers no lower-bound arithmetic for interruption reaction.
Our paper is exactly the missing latency-bound layer: it keeps revocability as the substrate
and asks how fast, per architecture, the agent can \emph{fully stop} while a tool waits.

\subsection{Endpointing, barge-in, and interruption policy}
Endpointing is the known latency-versus-reliability trade of detecting when the user stopped
talking: firing late adds silence to the onset floor, firing early risks cutting the user off
\citep{raux2008optimizing}.
Subsequent deployed work extracts the barge-in vs response-delay trade at scale
\citep{raux2012optimizing}.
Bohus and Horvitz \citep{bohus2011interruptions} go further and treat interruption as an
engagement \emph{decision}: sometimes the right behaviour is to keep talking, and deciding
whether to yield is policy, not physics.
These are the empirical and policy pillars our decomposition builds on.
Their gap is the same one: none states a lower bound on how long the agent's audio tail keeps
sounding after the user interrupts, nor how that bound depends on the architecture.

\subsection{Turn-taking timing and the two directions of responsiveness}
Human conversation shows modal gaps near $0$--$200$~ms with frequent sub-$100$~ms transitions
and ordered overlaps, not errors \citep{levinson2015timing,stivers2009universals,schegloff2000overlap}.
Quantitative corpus work further shows that a large share of transitions occur at or below
approximately $100$~ms, so a single threshold under-describes the distribution
\citep{heldner2010gaps}.
Backchannel placement is itself a timed behaviour: listeners tend to respond around a fixed
window after prosodic cues \citep{ward2000backchannel}.
Synthetic agents, for their part, lack a standard method for evaluating turn-taking behaviour
at all \citep{skantze2021turn}.
We use this literature to separate two responsiveness directions: producing one's first word
(onset) and releasing the floor (cessation), and we treat overlap as an orderly behaviour the
agent must be fast enough to support.

\subsection{Native full-duplex models}
End-to-end speech models such as Moshi demonstrate full-duplex dialogue with roughly $200$~ms of
streaming latency and visible interruption handling \citep{defossez2024moshi}, and commercial
APIs expose native speech-in/speech-out with built-in function calling
\citep{openai2024realtimeapi,google2024geminivoice}.
These systems are the architectural contrast for our per-topology cessation analysis: their
interruption reaction is bounded by codec-frame granularity, whereas a cascade's is bounded by
whatever audio is already in its synthesis queue.

\subsection{Agent orchestration, tools, and inference}
ReAct interleaves reasoning and tool acting \citep{yao2023react}, and graph frameworks such as
LangGraph make the resulting loops stateful and check-pointed \citep{langgraph2024}.
Chain-of-thought prompts make each reasoning step a visible latency component
\citep{wei2023chainofthought}.
Speculative decoding accelerates the reducible decode term without touching first-token or tool
delay \citep{leviathan2023fast}, and streaming long-context service is a serving-layer
concern \citep{xiao2023streamingllm}.
Finally, guardrail classification is increasingly placed inside the pipeline, where it must fit
a deadline rather than run after the fact \citep{dong2024building}.
This body of work motivates our setting: agents that reason, call tools, and are constrained by
a latency budget.
What it lacks is a formal treatment of \emph{interruptibility} as a schedulable property of
the orchestration rather than a property of any single component.

A separate line of engineering practice attacks tool latency directly: retrieval
round-trips are overlapped with drafting \citep{specrag2024}, intent prefetching lets an
assistant answer without a tool call at all \citep{agentsnotenough2024}, and the newest
full-duplex benchmarks evaluate tool use under streaming conditions
\citep{lin2026fdbv3}.
These are empirical existence proofs of the type-B/C hiding window we formalise in
Theorem~\ref{thm:tool_hiding}; none of them states a composition rule for when hiding is
possible or what unhideable type-A cost remains.
On the orchestration-operational side, graph frameworks already expose interrupt and
checkpoint primitives for human-in-the-loop reentry \citep{langgraph2024}; our point is that
such primitives are treated as operational features rather than as a schedulable latency
property.

\subsection{Recent full-duplex and interruption benchmarks}
A concurrent 2025--26 wave of benchmarks measures turn-taking and interruption behaviour
empirically.
Full-Duplex-Bench evaluates overlap handling in full-duplex spoken dialogue models
\citep{lin2025fullduplexbench}, its v3 extends this to tool use under streaming conditions
\citep{lin2026fdbv3}, and Easy Turn addresses turn-taking detection robustness
\citep{li2025easyturn}.
IHBench evaluates post-interruption recovery in structured-workflow voice agents
\citep{salimi2026ihbench}.
These works define and measure interruption and recovery \emph{outcomes} -- task success,
overlap accuracy, recovery outcome -- empirically.
None of them derives a formal lower bound on reaction latency per architecture (cascade vs
full-duplex vs decoupled), nor the phase-conditioned decomposition of the stop path.
The present paper supplies that missing formal account and is thus complementary to, not in
competition with, this benchmark wave.

\subsection{Position relative to the onset floor}
Our companion paper \citep{latencyfloor2025} bounds response onset for cascades and derives a
$k$-step silence-gap bound
for tool-using agents.
It treats the orchestration inside the language stage as a black box and stops at onset,
closing its open-problems discussion by explicitly flagging barge-in reaction as
untreated.
This paper opens that box: it bounds the other direction of responsiveness and shows that the
orchestrator's scheduling policy, not the components, decides whether reducible latency is
hidden or paid.
Its reactive cessation floor and the companion's predictive turn-end ceiling together delimit
the intervention window an agent can exploit.
